\section{Universality and the Heroic Life}

\begin{quotex}
Between 1931 and 1943 Julius Evola published 135 articles in the journal La Vita Italiana. What follows is the conclusion to ``Universalità imperiale e particolarismo nazionalistico'' from the April 1931 issue, in which Evola contrasts the Universality of the Empire to the particularity of the nation. In a few short paragraphs, he brings in the expected themes of caste, spirit, the ascetic life, the heroic life, fate, and so on.
\end{quotex}


\paragraph{Universality as knowledge and universality as action: the two bases of every imperial age.}


Knowledge is universal when it reaches the point of giving us the meaning of things, in the face of that greatness and eternity, everything that is \emph{pathos} and the fads of men vanishes: when it leads us into the primordial, the cosmic, into that which in the field of the spirit has the same characteristics of the purity and power of the oceans, the deserts, the ice fields. Every true universal tradition carried in itself this breath of vastness, animating disinterested forms of activity by it, arousing the sensibility for values that do not let themselves be measured by any criterion of utility and passion, whether individual or collective; introducing to the ``living'' a ``more than living''. This is the type of an invisible empire that history shows us in the examples of Brahmanic India, the Catholic Middle Ages, even of Hellenism: a homogeneous culture that rules every politically and economically conditioned reality, from the inside in a range independent of the people or city.

But we can think of a concept of empire, visible as well as invisible, having a both a material and a spiritual unity. We understand a similar empire when along with universality as knowledge, we also have \emph{universality as action}. Here as historical references, we can mention ancient China, Rome (in part), and once again the Middle Ages in the movements of the Crusades on the one hand and Islam on the other.

Universalizing action is pure action: it is \textbf{heroism}. So in the two conditions of imperiality we again find exactly the qualities that defined the two higher castes of antiquity, the sapiential caste (which is not necessarily the sacerdotal) and the warrior caste. We immediately point out that the concept of ``heroism'' here described is not that of the moderns. In the traditional concept, heroism is an ascesis in the strictest meaning of the word, and the hero is a nature so purified from the ``human'' elements, as far as he is the ascetic: he participates in the same character of purity as the great forces of circumstances, and has little to do with the passionality, sentimentality, and various movements of men, whether ideal or material, individual or collective.

The specific functions of each one of the ancient castes expressed its proper nature, the way of being of those who belonged to it: so war to the warrior was valued as his purpose, as the life for his own spiritual realization. He fought therefore in a ``pure'' way, war in itself was a good, and heroism was a ``pure'', hence universal, form, of activity. The rhetoric of the ``battle for rights'', ``territorial claims'', sentimental or humanitarian pretexts, and so on, are modern things in everything and for everyone, totally alien to the traditional concept of heroism. In the \emph{Bhagavad Gita}, in the \emph{Koran}, in the Latin concept of the \emph{mors triumphalis}, in the Hellenic assimilation of the hero at the beginning of the city, in the Nordic symbol of Valhalla, opened up only for the heroes, in certain aspects of the ``holy war'' known even by Catholic feudalism: we find, formulated in various ways, the transcendent idea, both supernatural and super-human, of heroism. Heroism here is a method of virile ascesis, of the conquering of one's lower nature, a life of immortality and in relationship with the eternal. Transfigured in a similar atmosphere, action gains a universal nature: it becomes almost a force from above, also capable of transposing the universality of a tradition of the spirit into a earthly body: it is the condition of the empire, in its highest meaning.

Are these exhumations of ancient ideals anachronistic and vain? They can also be so. But then this means only that current conditions are such to reduce the evocation of ideals and symbols, that today have lost their original meaning, to an empty rhetoric in which many indulge.

This does not block anything in the doctrinal sense. One can and must always trace a line of demarcation between concept and concept and to take care not to allow in a contradiction. When the points of reference are ``national pride'', the ``irredentistic demand'', the ``necessity for expansion'', etc., one is in the legitimate principles of a strongly modern nation, but one is certainly not in that of the empire. One wonders if a Roman had ever fought for any similar reason, and if he had ever needed to get motivated by such passional rhetoric to accomplish the miracle of that worldwide conquest, through which the universality of the luminous Greco-Roman civilization was spread to distant parts of the earth.

It is necessary to go back to the state of pure force, of force that comes with the same fatality and the same purity and the same inhumanity of the great forces of circumstances. The great conquerors always felt almost like \textbf{sons of fate}, carriers of a force that had to be realized and to which, starting from their own person, their own pleasure, their own tranquility, all of which had to be subdued and sacrificed. In its integral meaning, the empire is something higher, transcendent: the \emph{Sacrum Imperium} Holy Empire . How can we then associate the myth of the empire to this or that idealism or traditionalism (in a narrow sense) or sentimentalism or utilitarianism? How to connect it to the needs of a faction or nation, not to also mention a district, village, or country? Nevertheless, among the moderns it is all too frequently the case that it ends up in similar absurdities.

Whoever recalls the imperial symbols, which is the land that gave life to his body, must be capable of seeing all that. It is necessary that he knows what the ``nation'' is and what the ``empire'' is: which is the limit of one thing and that of the other. It is necessary that the mind be opened up to that which in man neither begins nor ends in man: that he understands, as the culmination of the most intense individuality, universality, both as knowledge and as action. And above all it is necessary that having the sense of the measure to which today everything is unnaturally reduced, he knows that there is all one world to which to say ``no'' before the auroral brightness of a possible European imperial era beyond the world of slaves and merchants.


\flrightit{Posted on 2011-08-10 by Aeneas}

\begin{center}* * *\end{center}

\begin{footnotesize}
\begin{sffamily}
\texttt{logres on 2011-08-12 at 21:45 said:}

``... which does not begin nor end with man... ''

Christianity has forgotten that the Father is archetypal and hence beyond our current form, and that our destiny is not to be merely man, but to pass to another state. How shameful it is to have to be reminded of this from a ``pagan'' -- but how instructive!


\hfill

\end{sffamily}
\end{footnotesize}

